\documentclass{article}
\usepackage[utf8]{inputenc}
\usepackage{geometry}
\usepackage{listings}
\usepackage{xcolor}
\usepackage{hyperref}
\usepackage{parskip}

\geometry{a4paper, margin=1in}

\definecolor{codebg}{RGB}{240, 240, 240}
\lstset{
    backgroundcolor=\color{codebg},
    basicstyle=\ttfamily\footnotesize,
    breaklines=true,
    frame=single,
    numbers=left,
    numberstyle=\tiny\color{gray},
    captionpos=b,
    tabsize=2
}

\title{React Native Expo Build Instructions}
\author{}
\date{\today}

\begin{document}

\maketitle

\section{Introduction}
This document provides complete instructions for generating APK and AAB files starting from a fresh clone of a React Native Expo project. The instructions cover both Linux/macOS and Windows operating systems.

\section{Prerequisites}
\begin{itemize}
    \item Node.js (v16 or later recommended)
    \item Java Development Kit (JDK) 11-17
    \item Android Studio with:
    \begin{itemize}
        \item Android SDK (API level 33+)
        \item Android Emulator
        \item Build Tools (33.0.0+)
    \end{itemize}
    \item Git installed
    \item Python (for Windows users)
\end{itemize}

\section{Initial Setup}

\subsection{Clone the Project}
\begin{lstlisting}[language=bash, caption=Clone the repository]
git clone <repository-url>
cd <project-folder>
\end{lstlisting}

\subsection{Install Dependencies}
\begin{lstlisting}[language=bash, caption=Install dependencies]
npm install
# or if using yarn
yarn install
\end{lstlisting}

\section{Build Preparation}

\subsection{Generate Android/iOS Directories (Expo)}
\begin{lstlisting}[language=bash, caption=Generate native directories]
npx expo prebuild
# or for older Expo versions
expo prebuild
\end{lstlisting}

\subsection{For Windows Users: Environment Setup}
\begin{itemize}
    \item Install Python 2.7 (required for some native modules)
    \item Set environment variables:
    \begin{itemize}
        \item ANDROID\_HOME to your SDK path (typically C:\textbackslash Users\textbackslash [username]\textbackslash AppData\textbackslash Local\textbackslash Android\textbackslash Sdk)
        \item Add platform-tools to PATH
    \end{itemize}
\end{itemize}

\section{Generating the Bundle}

\subsection{Linux/macOS}
\begin{lstlisting}[language=bash, caption=Linux/macOS bundle generation]
# Create assets directory
mkdir -p android/app/src/main/assets

# Generate bundle
npx react-native bundle \
    --platform android \
    --dev false \
    --entry-file index.js \
    --bundle-output android/app/src/main/assets/index.android.bundle \
    --assets-dest android/app/src/main/res/
\end{lstlisting}

\subsection{Windows}
\begin{lstlisting}[language=bash, caption=Windows bundle generation]
# Create assets directory
mkdir android\app\src\main\assets

# Generate bundle
npx react-native bundle ^
    --platform android ^
    --dev false ^
    --entry-file index.js ^
    --bundle-output android\app\src\main\assets\index.android.bundle ^
    --assets-dest android\app\src\main\res/
\end{lstlisting}

\section{Building APK/AAB}

\subsection{For Both Platforms}
\begin{lstlisting}[language=bash, caption=Build commands]
# Navigate to android directory
cd android

# Clean build (recommended)
./gradlew clean  # Linux/macOS
gradlew clean    # Windows

# Build APK
./gradlew assembleRelease  # Linux/macOS
gradlew assembleRelease    # Windows

# Build AAB (for Play Store)
./gradlew bundleRelease    # Linux/macOS
gradlew bundleRelease      # Windows
\end{lstlisting}

\section{Locating Build Outputs}
\begin{itemize}
    \item APK: \texttt{android/app/build/outputs/apk/release/app-release.apk}
    \item AAB: \texttt{android/app/build/outputs/bundle/release/app-release.aab}
\end{itemize}

\section{Troubleshooting}
\begin{itemize}
    \item \textbf{Missing assets}: Ensure bundle generation completed successfully
    \item \textbf{Java version issues}: Verify JDK 11-17 is installed
    \item \textbf{Windows path issues}: Use double backslashes or forward slashes in paths
    \item \textbf{Memory errors}: Add to \texttt{gradle.properties}:
    \begin{lstlisting}
org.gradle.jvmargs=-Xmx4096m -XX:MaxPermSize=1024m
    \end{lstlisting}
\end{itemize}

\section{Next Steps}
\begin{itemize}
    \item Test the APK on a device or emulator
    \item For production releases:
    \begin{itemize}
        \item Configure code signing
        \item Update version codes in \texttt{android/app/build.gradle}
        \item Consider using GitHub Actions or CI/CD for automated builds
    \end{itemize}
\end{itemize}

\end{document}